\documentclass[../../../include/open-logic-section]{subfiles}

\begin{document}

\olfileid{sth}{z}{story}
\olsection{కథ మరింత వివరంగా}
	
\olref[story][approach]{sec}లో
సమితుల నిర్మాణ ప్రక్రియ గురించి
షోన్‌ఫీల్డ్ వర్ణనను ఉటంకించాం.
ఇప్పుడు ఈ కథను కొంచెం
కచ్చితంగా చెప్పడానికి కొన్ని
సూత్రాలను రాయాలనుకుంటున్నాం.
అవి ఇవి:
\begin{enumerate}
\item[] \stageshier. ప్రతి సమితీ ఏదో ఒక దశలో ఏర్పడుతుంది.
\item[] \stagesord. దశలు క్రమంలో ఉంటాయి: కొన్ని ఇతర దశలకు \emph{ముందు}
ఉంటాయి.\footnote{దశలు \emph{సుక్రమితంగా} ఉన్నాయని
వాస్తవానికి మౌనంగా ఊహిస్తాం.
దీని అర్థం \olref[ordinals][]{chap}లో
వివరిస్తాం. ఇది బలమైన
ఊహ. నిజానికి \citet{Scott1974}
అందించిన తెలివైన పద్ధతితో
ఈ ఊహను \emph{తప్పించి}, తరువాత
\emph{వ్యుత్పాదించవచ్చు}.
(ప్రారంభ దశ ఉందని ఎందుకు
భావించాలో కూడా దీని ద్వారా
తెలుస్తుంది.) ఇక్కడ దానిలోకి
వెళ్లలేం; మరింత తెలుసుకోవడానికి
\citet{ButtonLT1} చూడండి.} 
\item[] \stagesacc. ఏ దశ $S$కైనా, $S$ దశకు
\emph{ముందు} ఏర్పడిన ఏ
సమితులనైనా తీసుకుంటే:
ఖచ్చితంగా ఆ సమితులే
మూలకాలుగా గల ఒక సమితి
$S$ దశలో ఏర్పడుతుంది.
$S$ దశలో మరేదీ ఏర్పడదు.
\end{enumerate}
ఇవి అనౌపచారిక సూత్రాలు;
అయినా వీటితో సెర్మెలో సమితి
సిద్ధాంతంలోని పలు స్వీకృతాలను
సమర్థించగలుగుతాం.

(ఒక జాగ్రత్త చెప్పాలి.
మనం పూర్తిగా ప్రామాణిక
స్వీకృతాలను, వాటి ప్రామాణిక
పేర్లతో పరిచయం చేస్తాం.
కానీ ఇప్పుడే ఇచ్చిన ఇటాలిక్
అక్షరాల్లోని సూత్రాలకు
సాహిత్యంలో ప్రత్యేక పేర్లు లేవు.
ఉపయోగకరంగా ఉంటాయని ఆశించి
మేమే వాటికి పేర్లు పెట్టాం.)

\end{document}
